\section{Topology-Aware Formulas}

This section documents the new default formulas used when
\texttt{serving\_scale} provides a model, cluster, interconnect, and topology.
The implementation lives in
\texttt{ServingCostModel::estimate\_stage\_breakdown(...)} in
\texttt{astra-sim/workload/ServingCostModel.cc}, with communication primitives
from \texttt{astra-sim/analytical/AnalyticalCostModels.cc}.

\subsection{Stage Decomposition}

For each scheduled prefill or decode batch, the simulator computes:

\begin{equation*}
T_{\mathrm{stage}} =
T_{\mathrm{base}} +
T_{\mathrm{attn}} +
T_{\mathrm{ffn/expert}} +
T_{\mathrm{tp}} +
T_{\mathrm{pp}} +
T_{\mathrm{ep}} +
T_{\mathrm{dpattn}} +
T_{\mathrm{pd}}
\end{equation*}

The stage multiplier is:

\begin{equation*}
m_{\mathrm{prefill}} = \eta_p
\end{equation*}

\begin{equation*}
m_{\mathrm{decode}} = \eta_d \phi_d
\end{equation*}

where \(\eta_p\) and \(\eta_d\) are the configured
\texttt{batch\_efficiency} terms and \(\phi_d\) is the decode
\texttt{interference\_factor}.

\subsection{Communication Micro-Models}

Let

\begin{equation*}
\alpha = \frac{L_{\mathrm{link}}}{10^9}
\end{equation*}

\begin{equation*}
\beta = \frac{1}{BW_{\mathrm{link}} \epsilon_{\mathrm{link}}}
\end{equation*}

The helper formulas reused from \texttt{AnalyticalCostModels.cc} are:

\begin{equation*}
\mathrm{AR}_{\mathrm{ring}}(B, n) =
2 (n - 1) \alpha + 2 \frac{n - 1}{n} B \beta
\end{equation*}

\begin{equation*}
\mathrm{AG}_{\mathrm{ring}}(B, n) =
(n - 1) \alpha + \frac{n - 1}{n} B \beta
\end{equation*}

\begin{equation*}
\mathrm{A2A}(B, n) =
(n - 1) \alpha + \frac{n - 1}{n} B \beta
\end{equation*}

\begin{equation*}
\mathrm{P2P}(B) =
\alpha + B \beta
\end{equation*}

\subsection{Attention and FFN or Expert Compute}

Let \(N_{\mathrm{tok}}\) be the sum of batch tokens, \(L\) the total number of
model layers, and \(p\) the pipeline-parallel degree. The simulator first
computes layers per stage as:

\begin{equation*}
L_{\mathrm{stage}} = \left\lceil \frac{L}{p} \right\rceil
\end{equation*}

Let \(S\) be the sequence-length hint used for the batch. The attention FLOPs
per token are:

\begin{equation*}
F_{\mathrm{attn/token}} = 8 H^2 + 4 S H
\end{equation*}

The default attention compute term is:

\begin{equation*}
\widehat{T}_{\mathrm{attn}} =
\frac{N_{\mathrm{tok}} L_{\mathrm{stage}} F_{\mathrm{attn/token}}}
{F_{\mathrm{peak}} \, tp}
\end{equation*}

For dense models, the FFN FLOPs per token are:

\begin{equation*}
F_{\mathrm{ffn/token}}^{\mathrm{dense}} = 6 H F
\end{equation*}

For MoE models, the default router plus expert FLOPs per token are:

\begin{equation*}
F_{\mathrm{ffn/token}}^{\mathrm{moe}} =
2 H E + 6 A H H_e
\end{equation*}

The FFN or expert compute term then becomes:

\begin{equation*}
\widehat{T}_{\mathrm{ffn/expert}} =
\frac{N_{\mathrm{tok}} L_{\mathrm{stage}} F_{\mathrm{ffn/token}}}
{F_{\mathrm{peak}} \, p_{\mathrm{ffn}}}
\end{equation*}

with

\begin{equation*}
p_{\mathrm{ffn}} =
\begin{cases}
tp & \text{dense} \\
tp \cdot ep & \text{MoE}
\end{cases}
\end{equation*}

\subsection{TP, PP, EP, and DP-Attention Terms}

Let \(b_a\) be the activation byte width used by the code path:

\begin{equation*}
b_a =
\begin{cases}
\texttt{bytes\_per\_activation} & \text{if configured} \\
b_{\mathrm{kv}} & \text{otherwise}
\end{cases}
\end{equation*}

The tensor-parallel message size is:

\begin{equation*}
Q_{\mathrm{tp}} = N_{\mathrm{tok}} H b_a
\end{equation*}

The TP collective term is:

\begin{equation*}
\widehat{T}_{\mathrm{tp}} =
\begin{cases}
2 L_{\mathrm{stage}} \,
\mathrm{AR}_{\mathrm{ring}}(Q_{\mathrm{tp}}, tp) & tp > 1 \\
0 & tp = 1
\end{cases}
\end{equation*}

The PP activation-send plus bubble term is:

\begin{equation*}
\widehat{T}_{\mathrm{pp}} =
\begin{cases}
\mathrm{P2P}(Q_{\mathrm{tp}}) +
\dfrac{(pp - 1)
\left(
\widehat{T}_{\mathrm{attn}} +
\widehat{T}_{\mathrm{ffn/expert}} +
\widehat{T}_{\mathrm{tp}}
\right)}
{\max(1, N_{\mathrm{tok}})} & pp > 1 \\
0 & pp = 1
\end{cases}
\end{equation*}

For MoE, the expert-dispatch payload is:

\begin{equation*}
Q_{\mathrm{ep}} = 2 N_{\mathrm{tok}} H b_a L_{\mathrm{stage}}
\end{equation*}

and the EP term is:

\begin{equation*}
\widehat{T}_{\mathrm{ep}} =
\begin{cases}
\mathrm{A2A}(Q_{\mathrm{ep}}, ep) & ep > 1 \text{ and MoE} \\
0 & \text{otherwise}
\end{cases}
\end{equation*}

The DP-attention synchronization payload per shard is:

\begin{equation*}
Q_{\mathrm{dpattn}} =
\frac{N_{\mathrm{tok}} H_{\mathrm{kv}} D_h \cdot 2 b_{\mathrm{kv}} L_{\mathrm{stage}}}
{d_a}
\end{equation*}

and the synchronization term is:

\begin{equation*}
\widehat{T}_{\mathrm{dpattn}} =
\begin{cases}
\mathrm{AG}_{\mathrm{ring}}(Q_{\mathrm{dpattn}}, d_a) & d_a > 1 \\
0 & d_a = 1
\end{cases}
\end{equation*}

\subsection{Scaled Final Component Values}

Let \(s_c\) be \texttt{compute\_scale} and \(s_m\) be \texttt{comm\_scale}. If a
per-token component override is set in the config, the simulator uses the
configured value directly:

\begin{equation*}
T_x =
m_{\mathrm{stage}} s_x N_{\mathrm{tok}} c_x
\end{equation*}

where \(s_x = s_c\) for compute components and \(s_x = s_m\) for communication
components.

Otherwise, the analytical defaults above are used:

\begin{equation*}
T_{\mathrm{attn}} = m_{\mathrm{stage}} s_c \widehat{T}_{\mathrm{attn}}
\end{equation*}

\begin{equation*}
T_{\mathrm{ffn/expert}} =
m_{\mathrm{stage}} s_c \widehat{T}_{\mathrm{ffn/expert}}
\end{equation*}

\begin{equation*}
T_{\mathrm{tp}} = m_{\mathrm{stage}} s_m \widehat{T}_{\mathrm{tp}}
\end{equation*}

\begin{equation*}
T_{\mathrm{pp}} = m_{\mathrm{stage}} s_m \widehat{T}_{\mathrm{pp}}
\end{equation*}

\begin{equation*}
T_{\mathrm{ep}} = m_{\mathrm{stage}} s_m \widehat{T}_{\mathrm{ep}}
\end{equation*}

\begin{equation*}
T_{\mathrm{dpattn}} = m_{\mathrm{stage}} s_m \widehat{T}_{\mathrm{dpattn}}
\end{equation*}

\begin{equation*}
T_{\mathrm{base}} = s_c B_{\mathrm{stage}}
\end{equation*}

\subsection{PD Transfer Stage}

The dedicated PD transfer stage uses:

\begin{equation*}
T_{\mathrm{pd}} =
L_{\mathrm{pd}} +
\frac{Q_{\mathrm{KV}}}{BW_{\mathrm{pd}} \epsilon_{\mathrm{pd}}}
\end{equation*}

This is assembled by
\texttt{ServingCostModel::estimate\_transfer\_breakdown(...)} and contributes to
the transfer-stage breakdown rather than to the prefill or decode stage.
